\scanpage{37}
\mspage{\handwritten{5}}
% archive sheet number: 36
est couvrant pour la topologie fpqc. Puisque $\operatorname{pr}_1\mathbin{o} v$ est fidèlement plat et de présentation finie, (resp. fidèlement plat et localement \deleted{pd}\inserted{de} présentation finie), il en est de même de $p$ (EGA IV 2.7.1), et comme $X$ est de présentation finie (resp. localement de présentation finie) sur $S$, il en est de même de $Y$ (EGA IV 17.7.4).

Montrons que ii) entraîne i). Il suffit de montrer que $p$ est couvrant pour la topologie fppf; or $p$ est fidèlement plat par hypothèse, et localement de présentation finie, puisque $X$ et $Y$ sont localement de présentation finie sur $S$.

\uline{10.9.}- \uline{Dans la situation rappelée au début de (10.0), soient $j\handwritten{\in}I$, et $G_j$ un $S_j$-groupe de présentation finie. Pour que $G^o=(G_j\handwritten{\times}_{S_j}S)^o$ soit représentable, il faut et il suffit qu'il existe \deleted{\unsure{un}} $i\handwritten{\in}I$, $i\handwritten{\geqslant}j$, tel que $(G_i)^o=(G_j\handwritten{\times}_{S_j}S_i)^o$ soit représentable.}

La condition est suffisante (3.3).

Montrons que la condition est nécessaire. \deleted{Alors}\inserted{\handwritten{En effet,}} $G^o$ est de présentation finie sur $S$ et ouvert dans $G$ (3.9). D'après (10.0), il existe $i\handwritten{\in}I$, $i\handwritten{\geqslant}j$, et un sous-préschéma en groupes ouvert $G_i^o$ de $G_i$ tel que $G_i^o\handwritten{\times}_{S_i}S=G^o$. Le morphisme structural $G^o\handwritten{\longrightarrow}S$ est connexe (i.e. à fibres géométriquement connexes)(Exp. $\mathrm{VI}_A$ 2.1.1). Alors, (EGA IV 9.3.3 et 9.7.7), quitte à augmenter $i$, on peut supposer que le morphisme structural $G_i^o\handwritten{\longrightarrow}S$ est connexe, donc que l'espace sous-jacent à $G_i^o$ n'est autre que \uline{$\handwritten{(}G_i^{\deleted{o}}\handwritten{)^o}$} (3.10.1), donc que $G_i^o$ représente $(G_i)^o$.

\uline{10.10.} \uuline{Rappelons de\deleted{\unsure{y|v}}\inserted{\handwritten{u}}x cas particuliers très utiles de la situation énoncée en (10.0)}

a) Etant donné un point $x$ d'un préschéma $X$, on pose $S_o=\operatorname{Spec}\handwritten{\mathbb{Z}}$\handwritten{/}et on considère la famille $(S_i)_{i\handwritten{\in}I}$ filtrante décroissante des voisinages \deleted{\unsure{principaux}} ouverts affines de $x$; alors $S=\operatorname{Spec}\handwritten{\mathcal{O}}_{X,x}$. En particulier, si $x$ est le point générique d'un schéma intègre $X$, on trouve $S=\handwritten{\kappa}(x)$.

b) On pose $S_o=\operatorname{Spec}\handwritten{\mathbb{Z}}$, et on considère la famille $(\handwritten{\mathscr{A}}_i)_{i\handwritten{\in}I}$ préordonnée par inclusion des sous-$\handwritten{\mathbb{Z}}$-algèbres de type fini de l'anneau d'un schéma affine $S$.
